%=====================================================================
\chapter{The Sixteen-Week Schedule}\label{app:schedule}
%=====================================================================
\normalfigures

This appendix is the \textbf{source of record} for the teaching schedule. The
one-page \texttt{ECAD\_Fall2026\_Lecture\_Schedule.pdf} distributed to students
is generated from it, and the slide pipeline parses the table below. Change
this first; regenerate the one-page version afterwards.

Twenty-three chapters are taught across sixteen weeks, so seven weeks carry two
chapters. The pairings are not arbitrary. In every case the two chapters are
either one topic the printed book splits for length --- the catalogue and its
variants, the two halves of the back office --- or a pair in which the second
is unusable without the first: a basket with nothing to check out, a gateway
with no payment method chosen.

One pairing is deliberate in the other direction. \textbf{Chapter~16, Security,
has a week to itself}, although it would pair naturally with Chapter~15. Its
laboratory attacks the student's own shop and then hardens it, and that does
not fit in half a week alongside anything else.

\section{The Schedule}

% MACHINE-READ. The slide builder takes the FIRST longtable in this file and
% requires: five columns; a bare integer week; chapters as comma-separated
% digits or N--M ranges; literal --- for an empty assessment cell; and no
% chapter appearing in two weeks.
\rowcolors{2}{white}{lightbg}
% column 3 holds the bold head "Chapters", which needs 1.8 cm at 10.95 pt
\begin{longtable}{@{}L{1.1cm}L{3.2cm}L{1.8cm}L{4.4cm}L{3.9cm}@{}}
\toprule
\rowcolor{primary!15}
\textbf{Wk} & \textbf{Lecture topic} & \textbf{Chapters} & \textbf{Lab} &
\textbf{Assessment} \\
\midrule
\endhead
1 & What e-commerce changed; business and revenue models & 1, 2 & Read four
real shops; the arithmetic of one order & --- \\
2 & Auctions, portals and communities; planning a framework & 3, 4 & Three
shapes compared; your requirements and your decision & --- \\
3 & The stack, and your shop in one command & 5 & Compose up; verify, break and
reset the shop & --- \\
4 & Products, categories and the data model beneath them & 6 & 50{,}000
products, and the index that rescues the category page & --- \\
5 & Product variations, options and user uploads & 7 & Option combinations, and
what an upload field accepts & \textbf{Quiz 1 (Ch 1--5)} \\
6 & Themes, layout and the user experience & 8 & Page weight before and after,
measured in bytes and requests & \textbf{Assignment 1 set} \\
7 & The shopping basket, and where state lives & 9 & Session, cookie and
database baskets across a restart & --- \\
8 & The checkout and the order state machine & 10 & Midterm, then queries
counted per checkout step & \textbf{Midterm (Ch 1--9)} \\
9 & Shipping, tax and zones; discounts and vouchers & 11, 12 & Geo zones and
GST; a voucher applied in the wrong order & \textbf{Assignment 1 due} \\
10 & Web payment systems; taking payment for orders & 13, 14 & A sandbox
gateway, and what happens when it times out & \textbf{Project proposal} \\
11 & User account management & 15 & Password cost against logins per second &
--- \\
12 & Securing the shop & 16 & The OWASP checks against your own container,
before and after hardening & \textbf{Quiz 2 (Ch 10--15)} \\
13 & Administration: the dashboard, orders, customers, refunds & 17, 18 & A
refund through the order state machine & \textbf{Assignment 2 set} \\
14 & Extending the platform: your first module & 19 & A shipping method and a
dashboard widget, written from nothing & --- \\
15 & Deployment, maintenance and performance; SEO & 20, 21 & Load test with
cache on and off; crawlability measured & \textbf{Assignment 2 due} \\
16 & Social, legal and ethical issues; project defence & 22, 23 & Project
presentations & \textbf{Project due; presentations} \\
\bottomrule
\end{longtable}

\section{How the Assessment Adds Up}

HEC publishes the assessment instruments for this course --- sessional
examination, home assignments, quizzes, project, presentations, final
examination --- without weights. The weights below are the department's, and
are stated here so that a student can see the whole shape of the semester on
one page.

\rowcolors{2}{white}{lightbg}
\begin{longtable}{@{}L{4.6cm}L{2.0cm}L{8.4cm}@{}}
\toprule
\rowcolor{primary!15}
\textbf{Instrument} & \textbf{Weight} & \textbf{What it examines} \\
\midrule
\endhead
Two quizzes & 10\% & Weeks 5 and 12. Short, closed-book, on the chapters named
in the schedule \\
Two assignments & 15\% & Set in weeks 6 and 13. Both are done on your own shop
and submitted as a change to it plus a short write-up \\
Midterm examination & 20\% & Week 8, Chapters 1--9 \\
Semester project & 25\% & The shop itself, built across the whole semester
(\chref{project}), submitted in week 16 \\
Presentation and defence & 5\% & Week 16. Ten minutes, plus questions you are
expected to be able to answer about your own decisions \\
Final examination & 25\% & All chapters, with emphasis after the midterm \\
\bottomrule
\end{longtable}

\begin{alertbox}[The project is not a week-16 activity]
Twenty-five per cent of the marks are for a shop that is built in the
laboratory of every chapter from \chref{stack} onward. A student who treats the
laboratories as optional homework and intends to build the project in the last
fortnight will be building, in two weeks and without help, something the rest
of the class built in fourteen weeks with a lecturer in the room.

\smallskip
\chref{project} states what is required and how it is marked. Read it in
Week~3, not in Week~15.
\end{alertbox}
